\chapter{总结与展望}

本章对全文研究工作进行归纳，提炼课题完成情况与工程价值，并结合实验结果提出后续可深化的研究方向。

\section{主要工作与创新点}

围绕"基于移动平台的固定目标识别及瞄准系统设计"，本文完成了系统总体方案制定、硬件平台搭建、软件闭环控制实现以及整机调试验证等工作。

在控制架构层面，本文提出并实现了"像素误差APID外环+角速度前馈内环"的双闭环控制结构。外环以目标中心像素偏差为输入，内环直接叠加HiPNUC IMU的角速度及角加速度前馈项，在平台机动或振动工况下可提前修正控制量，有效压缩纯视觉闭环的相位滞后，实现对静态目标的稳定指向。这一紧耦合的视觉+IMU双闭环架构是本课题的核心创新。

在系统架构层面，构建了"Linux视觉节点+STM32H7云台主控节点+独立底盘主控节点"的三节点异构分布式架构。三节点通过标准化接口协同，各自在最适合的计算平台内独立闭环运行，既降低了单节点负担，也提升了系统容错性。这种设计使得单节点降级不影响其他节点继续执行各自的控制任务，显著增强了整体系统的可靠性。

在工程实现层面，本文在"千元级成本、5\,W功耗、500\,g重量"三重约束下，完成了可落地、可复现的系统工程设计。相较于国外2万美元以上的高成本专业方案，本文为本科教学与小型无人系统研究提供了低成本参考路径，对"视觉+IMU"轻量化瞄准技术的推广具有示范意义。

此外，在底盘控制方面实现了基于八路灰度传感器的自主循迹、固定圈速闭环与编码器PI速度调节，并集成了A*路径规划算法，使底盘子系统具备从"轨迹跟随"到"自主规划导航"的扩展能力，突破了固定跑道的局限。

\section{不足与改进方向}

受限于实验条件与开发周期，当前系统在多个方面仍有提升空间。首先在视觉感知层面，目标检测在强光突变、目标部分遮挡与复杂背景下的鲁棒性不足。当前采用的HOG/HSV方法泛化能力相对较弱，不如深度学习检测器那样鲁棒。这是制约系统在各种实际光照条件下表现的主要因素。

其次在控制参数层面，像素闭环控制参数（PID增益及前馈系数）目前依赖手工整定，尚不具备参数自适应能力，对不同载体机动特性的适应性有限。这意味着在面对新的平台或运动特性时，需要重新调试参数。

第三，系统的时延补偿机制仍不完善。图像传输与控制指令执行之间的时间差在高动态工况下仍会引起一定的指向滞后，这在快速机动时会逐渐累积，影响追踪精度。

最后，工程中存在部分互斥保护未完全覆盖的并发风险，需要在后续版本中补充高负载工况下的稳定性测试。同时，可考虑通过扩展多传感器融合与更高性能执行机构，在保持系统轻量化的前提下进一步提升动态工况下的瞄准精度。

\section{后续研究工作展望}

后续工作可围绕多个方向展开。在视觉检测升级方面，可以以轻量化深度学习模型（如MobileNet-SSD或YOLO-nano）替代现有的HSV分割方案，以提升强光、遮挡与杂乱背景下的检测鲁棒性。同时轻量化模型可以保持在嵌入式平台上可接受的计算延迟，避免引入过高的系统负载。

在传感器融合层面，事件相机因其微秒级时间分辨率与低运动模糊特性，是高动态平台瞄准的潜在传感器升级路径。未来可以考虑让事件相机替换或补充现有USB摄像头，以完成视觉链路的升级。

在自适应控制方面，建立基于平台运动状态的在线参数自整定机制具有重要意义。同时引入预测性时延补偿，可进一步压缩高动态工况下的瞄准滞后，提升系统在实际运动中的追踪性能。

最后，将当前系统的控制框架封装为标准化软件模块，支持向不同嵌入式平台迁移，可以显著降低后续系统集成成本，提升技术可复用性。这对于推动相关技术在多个应用场景中的落地具有重要价值。

\newpage\quad %加一页，使得参考文献第一页的页眉为 参考文献
